% from ega0-19-fr.tex
\oldpage[0\textsubscript{IV-1}]{89}

\begin{proof}
La nécessité de la condition résulte de (19.4.2). Pour voir qu'elle est
suffisante, notons d'abord qu'en considérant un système fondamental de
voisinages ouverts de $0$ dans $A$ formés d'idéaux
$\mathfrak J_\alpha$ contenus dans $\mathfrak J$, on peut, en vertu de
(19.4.2), se ramener au cas où $A$ est discret puisque
$B/\mathfrak J_\alpha B$ est un
$(A/\mathfrak J_\alpha)$-module formellement projectif (19.2.5)~; par
définition d'un anneau préadmissible, $\mathfrak J$ est alors
\emph{nilpotent}. Il suffit en outre de démontrer la proposition lorsque
$\mathfrak J^2=0$, car si $\mathfrak J^m=0$, on l'appliquera successivement
aux anneaux $A_k=A/\mathfrak J^k$ et $B_k=B/\mathfrak J^kB$ et à l'idéal
$\mathfrak J^{k-1}/\mathfrak J^k$ de $A_k$ pour $k=2,3,\ldots,m$, en
notant (19.2.5) que $B_k=B\otimes_A A_k$ est un $A_k$-module formellement
projectif. Appliquons le critère (19.4.5), (ii) en considérant un idéal ouvert
$\mathfrak R$ de $B$ et un $2$-cocycle de Hochschild $A$-bilinéaire
symétrique $f$ de $(B/\mathfrak R)\times(B/\mathfrak R)$ dans un
$(B/\mathfrak R)$-module $L$. Considérons d'abord le cas particulier où
$\mathfrak JL=0$, de sorte que $L$ peut aussi être considéré comme un
$(B_0/\mathfrak RB_0)$-module, et $f$ se factorise en
\[
  (B/\mathfrak R)\times(B/\mathfrak R)\to
  (B_0/\mathfrak RB_0)\times(B_0/\mathfrak RB_0)
  \xrightarrow{f_0}L,
\]
où $f_0$ est un $2$-cocycle de Hochschild bilinéaire symétrique. En vertu de
l'hypothèse, il y a donc un idéal ouvert
$\mathfrak R'\subset\mathfrak R$ dans $B$ et une application
$A_0$-linéaire $g_0:B_0/\mathfrak R'B_0\to L$ telle que
\[
  f_0(\varphi_0(x),\varphi_0(y))
  =xg_0(y)-g_0(xy)+g_0(x)y
\]
pour $x,y$ dans $B_0/\mathfrak R'B_0$, $\varphi_0$ étant
l'homomorphisme canonique
$B_0/\mathfrak R'B_0\to B_0/\mathfrak RB_0$. On en conclut aussitôt que
l'application $A$-linéaire composée
$g:B/\mathfrak R'\to B_0/\mathfrak R'B_0\xrightarrow{g_0}L$ est telle
que $dg=f\circ(\varphi\times\varphi)$, où
$\varphi:B/\mathfrak R'\to B/\mathfrak R$ est l'homomorphisme
canonique.

Passons maintenant au cas général, et considérons d'abord le
$(B/\mathfrak R)$-module $L/\mathfrak JL=L'$, pour lequel on a
$\mathfrak JL'=0$~; si $f'$ est l'application $A$-bilinéaire composée
$(B/\mathfrak R)\times(B/\mathfrak R)\xrightarrow{f}L\to L'$, $f'$ est
encore un $2$-cocycle de Hochschild symétrique, et en vertu de ce qui précède,
il existe un idéal ouvert $\mathfrak R'\subset\mathfrak R$ dans $B$ et
une application $A$-linéaire $g':B/\mathfrak R'\to L'$ vérifiant
$dg'=f'\circ(\varphi'\times\varphi')$ pour l'application canonique
$\varphi':B/\mathfrak R'\to B/\mathfrak R$. Comme $B$ est un
$A$-module formellement projectif, il existe un idéal ouvert
$\mathfrak R'_1\subset\mathfrak R'$ et une application $A$-linéaire
$g_1:B/\mathfrak R'_1\to L$ telle que l'homomorphisme
$B/\mathfrak R'_1\to B/\mathfrak R'\xrightarrow{g'}L'$ se factorise en
$B/\mathfrak R'_1\xrightarrow{g_1}L\to L'$. Considérons alors le
$2$-cocycle de Hochschild
\[
  f_1(x,y)=f(\varphi_1(x),\varphi_1(y))-xg_1(y)+g_1(xy)-g_1(x)y,
\]
application $A$-bilinéaire symétrique de
$(B/\mathfrak R'_1)\times(B/\mathfrak R'_1)$ dans $L$. Le choix de
$g_1$ entraîne que $f_1$ prend ses valeurs dans $\mathfrak JL$. Comme
$\mathfrak J(\mathfrak JL)=0$, on peut de nouveau appliquer le premier cas,
et il y a donc un idéal ouvert
$\mathfrak R''\subset\mathfrak R'_1$ et une application $A$-linéaire
$g_2:B/\mathfrak R''\to\mathfrak JL$ telle que
\[
  f_1(\varphi_2(x),\varphi_2(y))=xg_2(y)-g_2(xy)+g_2(x)y
\]
dans $(B/\mathfrak R'')\times(B/\mathfrak R'')$,
$\varphi_2:B/\mathfrak R''\to B/\mathfrak R'_1$ étant l'application
canonique~; l'application $A$-linéaire
$g=g_2+g_1\circ\varphi_2:B/\mathfrak R''\to L$ vérifie donc
$dg=f\circ(\varphi\times\varphi)$ pour l'application canonique
$\varphi:B/\mathfrak R''\to B/\mathfrak R$. C.Q.F.D.
\end{proof}

\subsection{Lissité formelle et anneaux gradués associés}
\label{subsection:0.19.5-fr}

\begin{env}[19.5.1]
\label{0.19.5.1-fr}
